% =====================================================================
%  TA-IND-01 - Ingenieria de Requerimientos (ISR-401)
%  Facultad de Ciencias de la Computacion - UTEQ
%  Estudiante: Jeanpierre Robinson Espinoza (1722574223)
%
%  ARCHIVO PRINCIPAL. Compilar con pdfLaTeX + Biber:
%      pdflatex main.tex
%      biber main
%      pdflatex main.tex
%      pdflatex main.tex
% =====================================================================
\documentclass[12pt,a4paper]{article}

\usepackage[T1]{fontenc}
\usepackage[utf8]{inputenc}
\usepackage{lmodern}
\usepackage[spanish,es-noshorthands]{babel}
\usepackage{geometry}
\geometry{margin=2.5cm}
\usepackage{microtype}
\usepackage{graphicx}
\usepackage{booktabs}
\usepackage{longtable}
\usepackage{array}
\usepackage{tabularx}
\usepackage{enumitem}
\usepackage{xcolor}
\usepackage{csquotes}
\usepackage{caption}
\usepackage[hidelinks]{hyperref}
\usepackage[backend=biber,style=ieee]{biblatex}
\addbibresource{referencias.bib}

\newcommand{\pendiente}[1][]{\textcolor{red}{\textbf{[PENDIENTE#1]}}}
\newcommand{\pd}{\textcolor{red}{\scriptsize\textbf{PEND}}}
\newcolumntype{L}[1]{>{\raggedright\arraybackslash}p{#1}}

\newcommand{\captura}[3]{%
  \begin{figure}[htbp]
    \centering
    \IfFileExists{figuras/#1}{%
      \includegraphics[width=0.82\linewidth]{figuras/#1}%
    }{%
      \fbox{\parbox[c][5.5cm][c]{0.82\linewidth}{%
        \centering \pendiente{: reemplazar por la captura propia
        \texttt{figuras/#1}}}}%
    }%
    \caption{#2}
    \label{#3}
  \end{figure}%
}

\captionsetup{font=small,labelfont=bf}
\setlength{\parindent}{0pt}
\setlength{\parskip}{0.45em}
\setlength{\emergencystretch}{2em}
\renewcommand{\arraystretch}{1.2}

\begin{document}

\begin{center}
  {\Large\bfseries Requisitos de un sistema conocido a partir de\\[2pt]
  ISO/IEC/IEEE 29148:2018}\\[10pt]
  {\large TA-IND-01 - Trabajo autónomo individual}\\[6pt]
  Ingeniería de Requerimientos (ISR-401) - Cuarto nivel - Paralelo A\\
  Facultad de Ciencias de la Computación - Universidad Técnica Estatal de Quevedo\\[6pt]
  Jeanpierre Robinson Espinoza\\
  Código: 1722574223\\
  7 de noviembre de 2026
\end{center}

\vspace{0.4em}
\hrule
\vspace{0.4em}

% =====================================================================
\section{Descripción del sistema}
\label{sec:sistema}

El sistema analizado es InDrive, una aplicación móvil de transporte que
conecta a clientes que necesitan trasladarse con conductores particulares.
Su propósito es facilitar la movilidad de los usuarios mediante la
solicitud de viajes: el cliente indica su ubicación y su destino, los
conductores cercanos reciben la solicitud y proponen un precio, y el
cliente decide con qué oferta viajar. Intervienen tres tipos de usuarios.
Los clientes solicitan viajes y pagan el servicio; los conductores reciben
las solicitudes, ofertan un precio y realizan el traslado; y los
administradores supervisan el funcionamiento de la plataforma. El contexto
de uso es móvil y en movimiento, con conectividad variable, y el sistema
trata datos personales sensibles como la ubicación en tiempo real, el
número de teléfono y la información de pago.

Las Figuras~\ref{fig:cap1}, \ref{fig:cap2} y \ref{fig:cap3} son capturas
propias de la aplicación y respaldan los requisitos de la
Sección~\ref{sec:requisitos}.

\captura{captura-01.jpeg}{Captura propia de InDrive tomada durante el uso
del servicio (1 de 3).}{fig:cap1}

\captura{captura-02.jpeg}{Captura propia de InDrive tomada durante el uso
del servicio (2 de 3).}{fig:cap2}

\captura{captura-03.jpeg}{Captura propia de InDrive tomada durante el uso
del servicio (3 de 3).}{fig:cap3}

% =====================================================================
\section{Resumen analítico de ISO/IEC/IEEE 29148:2018}
\label{sec:resumen}

% ATENCION: este apartado debe escribirlo el estudiante a partir de la
% lectura de las clausulas 1, 3 y 5.2. Se deja marcado para no inventar
% contenido que no pueda sustentarse en la norma.

Fuentes de este apartado: \cite{iso29148} (cláusulas 1, 3 y 5.2), con
apoyo en \cite{pohl2025} y \cite{swebok2024}.

La norma ISO/IEC/IEEE 29148:2018 explica cómo identificar y documentar los
requisitos de un sistema o programa. También indica que es importante
conocer las necesidades de los usuarios para desarrollar un sistema que
cumpla con sus expectativas.

Los requisitos deben estar escritos de forma clara, precisa y fácil de
entender, para evitar confusiones durante el desarrollo. Además, la norma
presenta procesos para analizar, definir y revisar los requisitos de un
proyecto.

En conclusión, esta norma ayuda a organizar los requisitos y mantenerlos
actualizados para desarrollar un sistema que funcione correctamente.

% =====================================================================
\section{Glosario de términos}
\label{sec:glosario}

En la Tabla~\ref{tab:glosario} se recogen diez términos de la cláusula 3 de
\cite{iso29148}, con una definición parafraseada y un ejemplo propio del
sistema analizado.

{\small
\setlength{\tabcolsep}{3pt}
\begin{longtable}{@{}L{2.9cm}L{6.2cm}L{5.2cm}@{}}
\caption{Términos de la cláusula 3 aplicados a InDrive.}
\label{tab:glosario}\\
\toprule
\textbf{Término} & \textbf{Definición parafraseada} &
\textbf{Ejemplo en el sistema} \\
\midrule
\endfirsthead
\multicolumn{3}{@{}l}{\footnotesize\itshape Continuación de la Tabla~\ref{tab:glosario}}\\
\toprule
\textbf{Término} & \textbf{Definición parafraseada} &
\textbf{Ejemplo en el sistema} \\
\midrule
\endhead
\bottomrule
\endfoot
requisito (3.1.19) & Enunciado que expresa una necesidad junto con las
condiciones y restricciones que la acompañan. & ``El cliente debe aceptar
una oferta de viaje, siempre que su cuenta esté activa.'' \\
interesado (3.1.28) & Persona u organización con derecho, participación o
interés en el sistema y en que sus características satisfagan sus
necesidades. & El cliente, el conductor, el administrador y la empresa que
opera la plataforma. \\
usuario (3.1.35) & Individuo o grupo que interactúa con el sistema o se
beneficia de él durante su utilización. & El pasajero que solicita el viaje
desde su teléfono. \\
atributo (3.1.2) & Propiedad propia de una entidad que puede distinguirse
de manera cuantitativa o cualitativa. & La calificación promedio del
conductor, por ejemplo 4,8 de 5. \\
condición (3.1.6) & Atributo medible estipulado en un requisito que indica
la circunstancia o el evento en que este se aplica. & ``Si se interrumpe la
conexión, el viaje en curso se conserva.'' \\
restricción (3.1.7) & Limitación impuesta desde fuera al sistema, a su
diseño o al proceso con que se desarrolla. & Usar únicamente la pasarela de
pago autorizada por la entidad bancaria. \\
verificación de requisitos (3.1.26) & Confirmación, mediante examen, de que
los requisitos están bien formados. & Revisar que el requisito de rapidez
tenga métrica y criterio de aceptación antes de aprobarlo. \\
validación de requisitos (3.1.25) & Confirmación de que los requisitos
definen el sistema correcto según lo que esperan los interesados. & Consultar
a los conductores si las ofertas de precio cubren su necesidad real. \\
trazabilidad de requisitos (3.1.23) & Identificación y documentación del
origen y de la asignación de cada requisito dentro del conjunto. & El
requisito de cifrado se deriva de la necesidad de proteger los datos
personales del cliente. \\
sistema de interés (3.1.32) & Sistema cuyo ciclo de vida se analiza en el
contexto del documento. & La aplicación InDrive para pasajeros y
conductores, sin incluir la red de pagos del banco. \\
\end{longtable}
}

% =====================================================================
\section{Requisitos del sistema}
\label{sec:requisitos}

La Tabla~\ref{tab:requisitos} presenta quince requisitos de InDrive. Se
cubren seis tipos de la cláusula 5.2: funcional, de desempeño, de interfaz
externa, de calidad, de usabilidad y restricción. Cinco requisitos de
calidad se clasifican con la característica y la subcaracterística
correspondientes de ISO/IEC 25010:2023 \cite{iso25010} y se cuantifican con
una métrica.

{\footnotesize
\setlength{\tabcolsep}{3pt}
\begin{longtable}{@{}L{0.9cm}L{4.3cm}L{1.8cm}L{2.2cm}L{2.0cm}L{2.3cm}@{}}
\caption{Requisitos de InDrive con tipo, clasificación de calidad, métrica y
criterio de aceptación.}
\label{tab:requisitos}\\
\toprule
\textbf{ID} & \textbf{Enunciado} & \textbf{Tipo (5.2)} &
\textbf{ISO/IEC 25010:2023} & \textbf{Métrica} &
\textbf{Criterio de aceptación} \\
\midrule
\endfirsthead
\multicolumn{6}{@{}l}{\footnotesize\itshape Continuación de la Tabla~\ref{tab:requisitos}}\\
\toprule
\textbf{ID} & \textbf{Enunciado} & \textbf{Tipo (5.2)} &
\textbf{ISO/IEC 25010:2023} & \textbf{Métrica} &
\textbf{Criterio de aceptación} \\
\midrule
\endhead
\bottomrule
\endfoot
R01 & El sistema debe registrar una cuenta nueva con nombre, correo
electrónico y número de teléfono. & Funcional & -- & Campos válidos por
registro & El 100\,\% de los registros con los tres campos válidos crea la
cuenta; un campo inválido impide el registro y muestra el error. \\
R02 & El sistema debe autenticar al usuario con correo y contraseña, y
cerrar la sesión a petición del usuario. & Funcional & -- & Intentos de
acceso no autorizado & En 100 intentos con credenciales incorrectas se
conceden 0 accesos. \\
R03 & El sistema debe mostrar la ubicación actual del cliente y permitirle
seleccionar el destino sobre el mapa. & Funcional & -- & Destinos
registrados & La coordenada de destino queda registrada y visible antes de
confirmar la solicitud. \\
R04 & El sistema debe listar los conductores disponibles en un radio de
5 km alrededor del cliente. & Funcional & -- & Conductores listados por
radio & En una prueba con 3 conductores a menos de 5 km, se listan los 3. \\
R05 & El sistema debe permitir al conductor enviar una oferta de precio
para una solicitud recibida. & Funcional & -- & Ofertas entregadas & El
100\,\% de las ofertas enviadas llega al cliente con monto e identificación
del conductor. \\
R06 & El sistema debe permitir al cliente aceptar una oferta y debe marcar
como rechazadas las demás ofertas de esa solicitud. & Funcional & -- &
Ofertas mal asignadas & Al aceptar una oferta, el viaje queda asignado a un
solo conductor y las demás pasan a rechazadas en el 100\,\% de los casos. \\
R07 & El sistema debe mostrar nombre, fotografía, calificación y placa del
conductor asignado antes de iniciar el viaje. & Funcional & -- & Datos
visibles & Los 4 datos se muestran en el 100\,\% de los viajes asignados. \\
R08 & El sistema debe permitir calificar el viaje de 1 a 5 estrellas y
guardar la calificación asociada al viaje. & Funcional & -- & Calificaciones
registradas & El 100\,\% de las calificaciones enviadas queda asociada al
viaje correspondiente. \\
R09 & El sistema debe mostrar al cliente las ofertas de los conductores en
menos de 5 s desde que se emiten. & Desempeño & Eficiencia de desempeño:
comportamiento temporal & Tiempo de visualización de la oferta & El
percentil 95 de 100 ofertas se muestra en 5 s o menos con red 4G. \\
R10 & El sistema debe estar disponible para solicitar viajes el 99\,\% del
tiempo mensual. & Calidad & Fiabilidad: disponibilidad & Porcentaje de
tiempo disponible en 30 días & Disponibilidad mayor o igual al 99\,\%
(como máximo 7,2 h de indisponibilidad al mes). \\
R11 & El sistema debe cifrar el tráfico con TLS 1.2 o superior y exigir
contraseña para iniciar sesión. & Calidad & Seguridad: confidencialidad y
autenticidad & Porcentaje de tráfico cifrado y sesiones sin autenticar &
100\,\% del tráfico cifrado y 0 sesiones activas sin autenticación. \\
R12 & El cliente debe poder solicitar un viaje en no más de 5
interacciones desde la pantalla de inicio. & Usabilidad & Capacidad de
interacción: operabilidad & Número de interacciones hasta confirmar & En 3
de 3 pruebas, un usuario nuevo confirma el viaje en 5 interacciones o
menos. \\
R13 & El sistema debe restaurar el estado del viaje en curso en menos de
10 s después de una interrupción de conexión. & Calidad & Fiabilidad:
capacidad de recuperación & Tiempo de restauración del estado & En 9 de 10
pruebas con corte de red de 30 s, el estado se restaura en 10 s o menos. \\
R14 & El sistema no debe transferir la ubicación exacta del cliente a
terceros distintos del conductor asignado. & Restricción & -- & Flujos de
datos auditados & 0 transferencias de ubicación a destinatarios no
autorizados en la auditoría de flujos. \\
R15 & El sistema debe registrar cada pago en la pasarela externa y emitir
un comprobante con identificador único en 5 s o menos. & Interfaz externa &
-- & Tiempo de confirmación del pago & El 95\,\% de los pagos se confirma
en 5 s o menos y el 100\,\% genera comprobante con identificador único. \\
\end{longtable}
}

% =====================================================================
\section{Autoevaluación y corrección de los requisitos}
\label{sec:autoevaluacion}

Se revisó cada requisito contra las nueve características de 5.2.5 y
contra los términos prohibidos de 5.2.7. Tres requisitos no superaron la
revisión y se reescribieron; la Tabla~\ref{tab:autoevaluacion} conserva la
versión anterior y la corregida.

{\small
\setlength{\tabcolsep}{3pt}
\begin{longtable}{@{}L{0.9cm}L{2.7cm}L{5.3cm}L{5.4cm}@{}}
\caption{Requisitos que no superaron la revisión: versión anterior y
versión corregida.}
\label{tab:autoevaluacion}\\
\toprule
\textbf{ID} & \textbf{Característica o término que falla} &
\textbf{Versión anterior} & \textbf{Versión corregida} \\
\midrule
\endfirsthead
\multicolumn{4}{@{}l}{\footnotesize\itshape Continuación de la Tabla~\ref{tab:autoevaluacion}}\\
\toprule
\textbf{ID} & \textbf{Característica o término que falla} &
\textbf{Versión anterior} & \textbf{Versión corregida} \\
\midrule
\endhead
\bottomrule
\endfoot
R11 & Verificable. ``Proteger'' es ambiguo y el enunciado une dos
requisitos. & ``La aplicación debe proteger los datos personales mediante
cifrado y solicitar una contraseña para iniciar sesión.'' & ``El sistema
debe cifrar el tráfico con TLS 1.2 o superior y exigir contraseña para
iniciar sesión.'' \\
R12 & No ambigua. ``Pasos principales'' no define qué se cuenta como paso.
& ``El cliente debe poder solicitar un viaje en un máximo de 5 pasos
principales.'' & ``El cliente debe confirmar un viaje en no más de 5
interacciones desde la pantalla de inicio.'' \\
R09 & Completa. No indica condiciones de medición ni el percentil. & ``La
aplicación debe mostrar las ofertas de los conductores en menos de 5
segundos después de recibirlas.'' & ``El sistema debe mostrar las ofertas
en 5 s o menos en el percentil 95 de 100 solicitudes con red 4G.'' \\
\end{longtable}
}

% =====================================================================
\section{Reflexión}
\label{sec:reflexion}

El requisito que más me costó fue definir los de calidad, porque tuve que
pensar en valores medibles para aspectos como la rapidez, la seguridad y la
disponibilidad. En los requisitos funcionales bastaba con describir lo que
la aplicación hace, pero en los de calidad debía decidir cuánto era
suficiente y cómo se iba a comprobar. Al principio escribí enunciados como
``proteger los datos personales'', que suenan correctos pero no se pueden
verificar: no dicen qué se mide ni cuándo se considera cumplido. Al
aplicar la característica de verificabilidad de 5.2.5 me di cuenta de que
tenía que separar el cifrado de la autenticación y agregar una métrica
concreta. Lo mismo pasó con ``5 pasos principales'', donde no quedaba claro
qué contaba como paso. Esa fue la parte más difícil: convertir una
intención razonable en una condición que otra persona pueda medir sin
preguntarme a mí qué quise decir.

% =====================================================================
\section{Declaración de uso de inteligencia artificial}
\label{sec:ia}

Se declara el uso de un asistente de inteligencia artificial como apoyo
para la estructura del documento en LaTeX, la transcripción de las
respuestas del estudiante a las tablas del informe y la corrección de
errores de compilación. Las respuestas que sirven de base a los requisitos,
al glosario y a la reflexión fueron aportadas por el estudiante. Lo
obtenido se verificó compilando el documento, revisando cada requisito
contra la norma y contrastando los términos del glosario con la cláusula 3
de ISO/IEC/IEEE 29148:2018.

% =====================================================================
\printbibliography[title={Referencias}]

\end{document}
